\chapter{1977年江苏苏州市初试卷}

\section{解答题}

\begin{problem}
\begin{enumerate}
\item
计算：\( 1\div\left(-\frac{1}{8}\right)^2\times\left(-\frac{1}{4}\right)^3+\left(-1\frac{1}{3}\right)^2\div\left(\frac{2}{3}\right)^2 \).

\item
已知：\(3:x=6:y\)，求 \(\frac{x}{y}\) 的值.

\item
若 \(\sqrt{(3-x)^2}=x-3\) 成立，求 \(x\) 的取值范围.

\item
求函数 \(y=\frac{1}{2x^2-5x-3}\) 的自变量 \(x\) 的取值范围.

\item
若 \(0.0054=5.4\times10^n\)，求 \(n\).

\item
解方程：\((2x-3)(x+4)=-5\).

\item
求 \(\sin(-1020^\circ)\) 的值.

\item
已知 \(\lg2=0.3010\)，\(\lg3=0.4771\)，求 \(\lg1800\) 的值.

\item
求 \(\displaystyle\lim_{n\to\infty}\frac{2n+1}{n}\) 的值.

\item
求过点 \((3,4)\) 并平行于直线 \(2x+y-5=0\) 的直线方程.

\item
极坐标方程 \(\rho=3\) 表示什么曲线，并且将它化成直角坐标方程.

\item
某种反射镜面是旋转抛物面，它的镜面口直径是 \(40\mathrm{~cm}\)，深度是 \(5\mathrm{~cm}\)（如图），求抛物线的直角坐标方程和焦点坐标.

\end{enumerate}

\bitmapfigure[max width=0.36\linewidth]{img/1977/jiangsu_suzhou_city/q1_fig1.png}
\end{problem}

\begin{solution}
\begin{enumerate}
\item 将带分数化为假分数，得
\[
1\div\left(-\frac{1}{8}\right)^2\times\left(-\frac{1}{4}\right)^3+\left(-\frac{4}{3}\right)^2\div\left(\frac{2}{3}\right)^2
=1\div\frac{1}{64}\times\left(-\frac{1}{64}\right)+\frac{16}{9}\div\frac{4}{9}
=-1+4=3.
\]

\item 由比例式得 \(\frac{3}{x}=\frac{6}{y}\)，所以 \(3y=6x\)，即 \(y=2x\). 因此 \(\frac{x}{y}=\frac{1}{2}\).

\item 因为 \(\sqrt{(3-x)^2}=\lvert 3-x\rvert\)，而等式右边必须满足 \(x-3\geqslant0\). 当 \(x\geqslant3\) 时，\(\lvert 3-x\rvert=x-3\)，等式成立；当 \(x<3\) 时，\(\lvert 3-x\rvert=3-x\ne x-3\). 所以 \(x\) 的取值范围是 \([3,+\infty)\).

\item 分母不能为零，而
\[
2x^2-5x-3=(2x+1)(x-3).
\]
因此 \(x\ne-\frac{1}{2}\) 且 \(x\ne3\)，自变量 \(x\) 的取值范围是 \(\left\{x\in\R\mid x\ne-\frac{1}{2},\ x\ne3\right\}\).

\item 由 \(5.4\ne0\)，得
\[
10^n=\frac{0.0054}{5.4}=0.001=10^{-3}.
\]
所以 \(n=-3\).

\item 原方程化为
\[
2x^2+5x-12=-5,
\qquad 2x^2+5x-7=0,
\qquad (2x+7)(x-1)=0.
\]
所以 \(x=1\) 或 \(x=-\frac{7}{2}\).

\item 利用正弦函数的周期性，\(-1020^\circ+3\times360^\circ=60^\circ\)，所以
\[
\sin(-1020^\circ)=\sin60^\circ=\frac{\sqrt{3}}{2}.
\]

\item 因为 \(1800=18\times100=2\times3^2\times10^2\)，所以
\[
\lg1800=\lg2+2\lg3+2=0.3010+2\times0.4771+2=3.2552.
\]

\item 当 \(n\ne0\) 时，
\[
\frac{2n+1}{n}=2+\frac{1}{n}.
\]
由于 \(\displaystyle\lim_{n\to\infty}\frac{1}{n}=0\)，所以
\[
\displaystyle\lim_{n\to\infty}\frac{2n+1}{n}=2.
\]

\item 与直线 \(2x+y-5=0\) 平行的直线可设为 \(2x+y+C=0\). 将点 \((3,4)\) 代入，得 \(2\times3+4+C=0\)，所以 \(C=-10\). 所求直线方程为 \(2x+y-10=0\).

\item 极坐标方程 \(\rho=3\) 表示以极点为圆心、半径为 \(3\) 的圆. 由 \(x=\rho\cos\theta\)，\(y=\rho\sin\theta\) 以及 \(x^2+y^2=\rho^2\)，得其直角坐标方程为 \(x^2+y^2=9\).

\item 取抛物线顶点为原点，开口方向为 \(x\) 轴正方向，则抛物线可设为 \(y^2=4px\). 镜面口直径为 \(40\)，所以口的两个端点的纵坐标为 \(20\) 和 \(-20\)；深度为 \(5\)，所以这两个端点的横坐标为 \(5\). 代入抛物线方程，得 \(20^2=4p\times5\)，即 \(p=20\). 因此抛物线的直角坐标方程为 \(y^2=80x\)，焦点坐标为 \((20,0)\).
\end{enumerate}
\end{solution}

\begin{problem}
化简：\( \left[a^{-\frac{3}{2}}b\left(ab^{-2}\right)^{-\frac{1}{2}}\right]^3 \)，并求当 \(a=\frac{1}{\sqrt{2}}\)，\(b=\frac{1}{\sqrt[3]{2}}\) 时的值.
\end{problem}

\begin{solution}
因为题中给定的 \(a\)，\(b\) 均为正数，可以使用正数幂的运算性质. 先化简括号内的式子：
\[
\left(ab^{-2}\right)^{-\frac{1}{2}}=a^{-\frac{1}{2}}b,
\qquad
a^{-\frac{3}{2}}b\left(ab^{-2}\right)^{-\frac{1}{2}}=a^{-2}b^2.
\]
所以原式为
\[
\left(a^{-2}b^2\right)^3=a^{-6}b^6=\left(\frac{b}{a}\right)^6.
\]
又因为 \(a=2^{-\frac{1}{2}}\)，\(b=2^{-\frac{1}{3}}\)，故
\[
\left(\frac{b}{a}\right)^6
=\left(2^{-\frac{1}{3}+\frac{1}{2}}\right)^6
=\left(2^{\frac{1}{6}}\right)^6=2.
\]
所以化简结果为 \(\left(\frac{b}{a}\right)^6\)，给定条件下的值为 \(2\).
\end{solution}

\begin{problem}
已知 \(\cos\theta=0.6\)，并且 \(\theta\) 是第四象限的角，求 \(\sin2\theta\) 的值.
\end{problem}

\begin{solution}
因为 \(\cos\theta=0.6=\frac{3}{5}\)，且 \(\theta\) 在第四象限，所以 \(\sin\theta<0\). 由平方关系得
\[
\sin\theta=-\sqrt{1-\cos^2\theta}=-\sqrt{1-\left(\frac{3}{5}\right)^2}=-\frac{4}{5}.
\]
于是
\[
\sin2\theta=2\sin\theta\cos\theta=2\times\left(-\frac{4}{5}\right)\times\frac{3}{5}=-\frac{24}{25}.
\]
\end{solution}

\begin{problem}
将 \(40\) 分成两个正数，使它们的平方和为最小，求这两个数.
\end{problem}

\begin{solution}
设两个正数分别为 \(x\) 和 \(40-x\)，则 \(0<x<40\)，它们的平方和为
\[
S=x^2+(40-x)^2=2x^2-80x+1600=2(x-20)^2+800.
\]
因为 \((x-20)^2\geqslant0\)，所以 \(S\geqslant800\)，当且仅当 \(x=20\) 时取等号. 此时 \(40-x=20\)，故这两个数都是 \(20\).
\end{solution}

\begin{problem}
已知 \(AB\) 两地相距 \(600\text{ 公里}\)，从 \(A\) 地开往 \(B\) 地的客车每小时比货车快 \(15\text{ 公里}\)，客车比货车少用 \(2\) 小时，求客车和货车的速度.
\end{problem}

\begin{solution}
设货车的速度为 \(v\mathrm{~km/h}\)，则客车的速度为 \((v+15)\mathrm{~km/h}\)，其中 \(v>0\). 由客车比货车少用 \(2\) 小时，得
\[
\frac{600}{v}-\frac{600}{v+15}=2.
\]
因为 \(v(v+15)>0\)，所以
\[
600\times15=2v(v+15),
\qquad v^2+15v-4500=0,
\qquad (v-60)(v+75)=0.
\]
由 \(v>0\)，得 \(v=60\). 因此货车的速度为 \(60\mathrm{~km/h}\)，客车的速度为 \(75\mathrm{~km/h}\).
\end{solution}

\begin{problem}
\(MT\) 是 \(\odot O\) 的切线，\(O\) 是圆心，\(T\) 是切点，\(MO\) 与 \(\odot O\) 交于 \(A\) 点，已知 \(MA\) 的长为 \(h\)，\(\odot O\) 的半径为 \(r\)，\(TP\) 为 \(MO\) 的垂线，垂足是 \(P\)，求 \(TP\) 的长.

\bitmapfigure[float=inline,max width=0.52\linewidth]{img/1977/jiangsu_suzhou_city/q6_fig1.png}
\end{problem}

\begin{solution}
由图中点的位置，\(A\) 在 \(M\)，\(O\) 之间，且 \(OA=r\)，所以 \(MO=MA+AO=h+r\). 半径垂直于切线，故 \(OT\perp MT\)，在直角三角形 \(MTO\) 中，
\[
MT^2=MO^2-OT^2=(h+r)^2-r^2=h(h+2r).
\]
又因为 \(TP\perp MO\)，用三角形 \(MTO\) 的面积表示式得
\[
\frac{1}{2}MT\cdot OT=\frac{1}{2}MO\cdot TP.
\]
代入 \(OT=r\)，\(MO=h+r\) 以及 \(MT=\sqrt{h(h+2r)}\)，得
\[
TP=\frac{MT\cdot OT}{MO}=\frac{r\sqrt{h(h+2r)}}{h+r}.
\]
\end{solution}

\section{附加题}

\begin{problem}
如果二次方程 \((m+1)x^2+(2m-1)x+(m-1)=0\) 无实数根，求证：二次方程 \((m-3)x^2-2(m+3)x-(m+5)=0\) 一定有两个不相等的实根.
\end{problem}

\begin{solution}
由第一个方程无实数根，得其判别式小于零，即
\[
\Delta_1=(2m-1)^2-4(m+1)(m-1)=4m^2-4m+1-4m^2+4=5-4m<0.
\]
所以 \(m>\frac{5}{4}\). 按题意第二个式子为二次方程，故 \(m-3\ne0\). 对第二个方程，其判别式为
\[
\Delta_2=[-2(m+3)]^2-4(m-3)[-(m+5)]=8(m^2+4m-3)=8\left((m+2)^2-7\right).
\]
由于 \(m>\frac{5}{4}\)，所以 \(m+2>\frac{13}{4}>\sqrt{7}\)，从而 \(\Delta_2>0\). 因此第二个二次方程的二次项系数不为零且判别式大于零，必有两个不相等的实根.

这里使用了题目将第二个式子规定为“二次方程”这一条件；若只假定第一个式子为二次方程而不要求 \(m\ne3\)，则 \(m=3\) 时第二个式子退化为一次方程，原结论不成立.
\end{solution}

\begin{problem}
\(\triangle ABC\) 的面积为 \(10\sqrt{3}\)，周长为 \(20\)，并且三个内角成等差数列，求各边的长.
\end{problem}

\begin{solution}
不妨设 \(\angle B\) 是三个内角中的中间项，并记 \(a=BC\)，\(b=CA\)，\(c=AB\). 因为三角形内角和为 \(180^\circ\)，且三个内角成等差数列，所以 \(\angle B=60^\circ\). 由面积公式得
\[
10\sqrt{3}=\frac{1}{2}ac\sin60^\circ=\frac{\sqrt{3}}{4}ac,
\]
所以 \(ac=40\). 由余弦定理，
\[
b^2=a^2+c^2-2ac\cos60^\circ=a^2+c^2-ac=(a+c)^2-3ac.
\]
令 \(s=a+c\)，由周长为 \(20\) 得 \(b=20-s\)，于是
\[
(20-s)^2=s^2-120,
\Longrightarrow 400-40s+s^2=s^2-120
\Longrightarrow 40s=520
\Longrightarrow s=13.
\]
因此 \(b=7\)，且 \(a\)，\(c\) 是方程
\[
t^2-13t+40=0,
\qquad (t-5)(t-8)=0
\]
的两个根，所以 \(\{a,c\}=\{5,8\}\). 各边的长为 \(5\)，\(7\)，\(8\). 验算得 \(\frac{1}{2}\times5\times8\times\sin60^\circ=10\sqrt{3}\)，且三边满足三角形不等式，故结果成立.
\end{solution}
